\documentclass{amsart}
% For dealing with references we use the comment environment
\usepackage{verbatim}
\newenvironment{reference}{\comment}{\endcomment}
%\newenvironment{reference}{}{}
\newenvironment{slogan}{\comment}{\endcomment}
\newenvironment{history}{\comment}{\endcomment}

% For commutative diagrams we use Xy-pic
\usepackage[all]{xy}

% We use 2cell for 2-commutative diagrams.
\xyoption{2cell}
\UseAllTwocells

% We use multicol for the list of chapters between chapters
\usepackage{multicol}

% This is generally recommended for better output
\usepackage{lmodern}
\usepackage[T1]{fontenc}

% For cross-file-references

% Package for hypertext links:
\usepackage{hyperref}

% For any local file, say "hello.tex" you want to link to please

% Theorem environments.
%
\theoremstyle{plain}
\newtheorem{theorem}[subsection]{Theorem}
\newtheorem{proposition}[subsection]{Proposition}
\newtheorem{lemma}[subsection]{Lemma}

\theoremstyle{definition}
\newtheorem{definition}[subsection]{Definition}
\newtheorem{example}[subsection]{Example}
\newtheorem{exercise}[subsection]{Exercise}
\newtheorem{situation}[subsection]{Situation}

\theoremstyle{remark}
\newtheorem{remark}[subsection]{Remark}
\newtheorem{remarks}[subsection]{Remarks}

\numberwithin{equation}{subsection}

% Macros
%
\def\lim{\mathop{\mathrm{lim}}\nolimits}
\def\colim{\mathop{\mathrm{colim}}\nolimits}
\def\Spec{\mathop{\mathrm{Spec}}}
\def\Hom{\mathop{\mathrm{Hom}}\nolimits}
\def\Ext{\mathop{\mathrm{Ext}}\nolimits}
\def\SheafHom{\mathop{\mathcal{H}\!\mathit{om}}\nolimits}
\def\SheafExt{\mathop{\mathcal{E}\!\mathit{xt}}\nolimits}
\def\Sch{\mathit{Sch}}
\def\Mor{\mathop{\mathrm{Mor}}\nolimits}
\def\Ob{\mathop{\mathrm{Ob}}\nolimits}
\def\Sh{\mathop{\mathit{Sh}}\nolimits}
\def\NL{\mathop{N\!L}\nolimits}
\def\CH{\mathop{\mathrm{CH}}\nolimits}
\def\proetale{{pro\text{-}\acute{e}tale}}
\def\etale{{\acute{e}tale}}
\def\QCoh{\mathit{QCoh}}
\def\Ker{\mathop{\mathrm{Ker}}}
\def\Im{\mathop{\mathrm{Im}}}
\def\Coker{\mathop{\mathrm{Coker}}}
\def\Coim{\mathop{\mathrm{Coim}}}

% Boxtimes
%
\DeclareMathSymbol{\boxtimes}{\mathbin}{AMSa}{"02}

%
% Macros for moduli stacks/spaces
%
\def\QCohstack{\mathcal{QC}\!\mathit{oh}}
\def\Cohstack{\mathcal{C}\!\mathit{oh}}
\def\Spacesstack{\mathcal{S}\!\mathit{paces}}
\def\Quotfunctor{\mathrm{Quot}}
\def\Hilbfunctor{\mathrm{Hilb}}
\def\Curvesstack{\mathcal{C}\!\mathit{urves}}
\def\Polarizedstack{\mathcal{P}\!\mathit{olarized}}
\def\Complexesstack{\mathcal{C}\!\mathit{omplexes}}
% \Pic is the operator that assigns to X its picard group, usage \Pic(X)
% \Picardstack_{X/B} denotes the Picard stack of X over B
% \Picardfunctor_{X/B} denotes the Picard functor of X over B
\def\Pic{\mathop{\mathrm{Pic}}\nolimits}
\def\Picardstack{\mathcal{P}\!\mathit{ic}}
\def\Picardfunctor{\mathrm{Pic}}
\def\Deformationcategory{\mathcal{D}\!\mathit{ef}}

\setlength{\emergencystretch}{3em}
\begin{document}
\title[Stacks Project, Tag 07RY]{480 --- Universal degree stratification for separated flat locally finitely presented quasi-finite morphisms with universally bounded fibres}
\maketitle
\noindent Copyright (C) 2005--2025 Johan de Jong. Stacks Project authors. Source: \href{https://stacks.math.columbia.edu/tag/07RY}{Tag 07RY}.

\noindent Editorial difficulty note (not author text). Difficulty H2: The maximal-degree locus must be shown open with finite locally free restriction by étale splitting and elimination of the residual component. Induction on the degree must also prove the universal property for arbitrary nonreduced test schemes, using reduction and universal closedness to recover finiteness..

\noindent Editorial provenance note (not author text). History: excerpt from revision \texttt{a0444\allowbreak 6e57e\allowbreak c1fbc\allowbreak 252a8\allowbreak 71afc\allowbreak ec775\allowbreak 2fb28\allowbreak 07b14}, more-morphisms.tex, lines 12760--12842. Reference presentation was adapted; original-author context and core support blocks were retained or added. The mathematical wording is unchanged. Licensed under GNU FDL 1.2 or later; no invariant sections or cover texts. See ../licenses/COPYING. Context blocks reproduce author definitions and supporting results; they are not additional counted problems.

\begin{lemma}
\label{lemma-stratify-flat-fp-lqf-universally-bounded}
Let $f : X \to S$ be flat, locally of finite presentation, separated,
locally quasi-finite with universally bounded fibres. Then there exist
closed subsets
$$
\emptyset = Z_{-1} \subset Z_0 \subset Z_1 \subset Z_2 \subset
\ldots \subset Z_n = S
$$
such that with $S_r = Z_r \setminus Z_{r - 1}$ the stratification
$S = \coprod_{r = 0, \ldots, n} S_r$ is characterized by the following
universal property: Given $g : T \to S$ the projection
$X \times_S T \to T$ is finite locally
free of degree $r$ if and only if $g(T) \subset S_r$ (set theoretically).
\end{lemma}

\begin{proof}
Let $n$ be an integer bounding the degree of the fibres of $X \to S$.
By Morphisms, Lemma \href{https://stacks.math.columbia.edu/tag/03J7}{lemma base change universally bounded (Tag 03J7)}
we see that any base change has degrees of fibres bounded by $n$ also.
In particular, all the integers $r$ that occur in the statement of the lemma
will be $\leq n$. We will prove the lemma by induction on $n$. The base
case is $n = 0$ which is obvious.

\medskip\noindent
We claim the set of points $s \in S$
with $\deg_{\kappa(s)}(X_s) = n$ is an open subset $S_n \subset S$
and that $X \times_S S_n \to S_n$ is finite locally free of degree $n$.
Namely, suppose that $s \in S$ is such a point. Choose an elementary
\'etale morphism $(U, u) \to (S, s)$ and a decomposition
$U \times_S X = W \amalg V$ as in
Lemma \href{https://stacks.math.columbia.edu/tag/02LP}{lemma etale splits off quasi finite part (Tag 02LP)}.
Since $V \to U$ is finite, flat, and locally of finite presentation,
we see that $V \to U$ is finite locally free, see
Morphisms, Lemma \href{https://stacks.math.columbia.edu/tag/02KB}{lemma finite flat (Tag 02KB)}.
After shrinking $U$ to a smaller neighbourhood of $u$
we may assume $V \to U$ is finite locally free of some degree $d$, see
Morphisms, Lemma \href{https://stacks.math.columbia.edu/tag/04MH}{lemma finite locally free (Tag 04MH)}.
As $u \mapsto s$ and $W_u = \emptyset$ we see that $d = n$. Since $n$
is the maximum degree of a fibre we see that $W = \emptyset$!
Thus $U \times_S X \to U$ is finite locally free of degree $n$.
By Descent, Lemma \href{https://stacks.math.columbia.edu/tag/02VO}{lemma descending property finite locally free (Tag 02VO)}
we conclude that $X \to S$ is finite locally free of degree $n$
over $\Im(U \to S)$ which is an open neighbourhood of $s$
(Morphisms, Lemma \href{https://stacks.math.columbia.edu/tag/03WT}{lemma etale open (Tag 03WT)}).
This proves the claim.

\medskip\noindent
Let $S' = S \setminus S_n$ endowed with the reduced induced scheme
structure and set $X' = X \times_S S'$. Note that the degrees of fibres
of $X' \to S'$ are universally bounded by $n - 1$. By induction we find a
stratification $S' = S_0 \amalg \ldots \amalg S_{n - 1}$ adapted
to the morphism $X' \to S'$. We claim that $S = \coprod_{r = 0, \ldots, n} S_r$
works for the morphism $X \to S$. Let $g : T \to S$ be a morphism of schemes
and assume that $X \times_S T \to T$ is finite locally free of degree $r$.
As remarked above this implies that $r \leq n$. If $r = n$, then it is
clear that $T \to S$ factors through $S_n$. If $r < n$, then
$g(T) \subset S' = S \setminus S_d$ (set theoretically) hence
$T_{red} \to S$ factors through $S'$, see
Schemes, Lemma \href{https://stacks.math.columbia.edu/tag/0356}{lemma map into reduction (Tag 0356)}.
Note that $X \times_S T_{red} \to T_{red}$ is also
finite locally free of degree $r$ as a base change.
By the universal property of the stratification
$S' = \coprod_{r = 0, \ldots, n - 1} S_r$ we see that $g(T) = g(T_{red})$
is contained in $S_r$.
Conversely, suppose that we have $g : T \to S$ such that
$g(T) \subset S_r$ (set theoretically).
If $r = n$, then $g$ factors through $S_n$ and
it is clear that $X \times_S T \to T$
is finite locally free of degree $n$ as a base change.
If $r < n$, then $X \times_S T \to T$ is a morphism which is
separated, flat, and locally of finite presentation, such that
the restriction to $T_{red}$ is finite locally free of degree $r$.
Since $T_{red} \to T$ is a universal homeomorphism, we conclude
that $X \times_S T_{red} \to X \times_S T$ is a universal homeomorphism
too and hence $X \times_S T \to T$ is universally closed (as this
is true for the finite morphism $X \times_S T_{red} \to T_{red}$).
It follows that $X \times_S T \to T$ is finite, for example by
Lemma \href{https://stacks.math.columbia.edu/tag/02LS}{lemma characterize finite (Tag 02LS)}. Then we can use
Morphisms, Lemma \href{https://stacks.math.columbia.edu/tag/02KB}{lemma finite flat (Tag 02KB)}
to see that $X \times_S T \to T$ is finite locally free.
Finally, the degree is $r$ as all the fibres have degree $r$.
\end{proof}
\end{document}
